\begin{pastexam}[label={px:gk3-097-ds-04}]{高考三級資訊處理 97 年 第四題}
（申論；25 分）

Ackermann's Function A(m,n)的定義如下：
\[
A(m, n) = \begin{cases}
n + 1 & \text{, if } \ m = 0 \\
A(m - 1, 1) & \text{, if } \ n = 0 \\
A(m - 1, A(m, n - 1)) & \text{, otherwise}
\end{cases}
\]
此函數的成長速度相當快，對於 m 和 n 是很小時亦然
\begin{itemize}[leftmargin=2.2em, nosep]
\item[㈠] 試寫一遞迴演算法（Recursive Algorithm）來計算此函數值。（15 分）
\item[㈡] 試求算出 A(2,2)的值。（需列出求算過程）（10 分）
\end{itemize}
\end{pastexam}

\begin{pastsoln}
答案：㈡ $A(2, 2) = 7$。

\textbf{㈠} 定義本身就是遞迴，三種情況照順序判斷即可（$m, n$ 都是非負整數）：
\begin{pseudo}{Ackermann 函數}
\Function{Ackermann}{$m, n$}
    \If{$m = 0$}
        \State \Return $n + 1$
    \ElsIf{$n = 0$}
        \State \Return \Call{Ackermann}{$m - 1, 1$}
    \Else
        \State $t \gets$ \Call{Ackermann}{$m, n - 1$} \Comment{先算內層}
        \State \Return \Call{Ackermann}{$m - 1, t$}
    \EndIf
\EndFunction
\end{pseudo}
\textbf{會停止的理由。} 把參數 $(m, n)$ 依字典順序比大小（先比 $m$，$m$ 相同再比 $n$）。 第 5 行
呼叫 $(m - 1, 1)$、第 8 行呼叫 $(m - 1, t)$，$m$ 都變小；第 7 行呼叫 $(m, n - 1)$，$m$ 不變、$n$
變小。 每次遞迴呼叫的參數都嚴格變小，而非負整數對的字典順序沒有無限遞減的序列，所以遞迴一定在
有限步內走到 $m = 0$ 的情況停下來。 $A(n, n)$ 長得比任何只用事先定好次數的迴圈算得出來的函數
都快；資結書第~13.3~節互斥集合的 $\alpha(n)$ 就是它的某種反函數。

\textbf{㈡} 由內往外算，每一行都只用定義的三條規則：
\begin{align*}
A(1, 0) &= A(0, 1) = 2 \\
A(1, 1) &= A(0, A(1, 0)) = A(0, 2) = 3 \\
A(2, 0) &= A(1, 1) = 3 \\
A(1, 2) &= A(0, A(1, 1)) = A(0, 3) = 4 \\
A(1, 3) &= A(0, A(1, 2)) = A(0, 4) = 5 \\
A(2, 1) &= A(1, A(2, 0)) = A(1, 3) = 5 \\
A(1, 4) &= A(0, A(1, 3)) = A(0, 5) = 6 \\
A(1, 5) &= A(0, A(1, 4)) = A(0, 6) = 7 \\
A(2, 2) &= A(1, A(2, 1)) = A(1, 5) = 7
\end{align*}
所以 $A(2, 2) = 7$。 照㈠的演算法實際執行，從 $A(2, 2)$ 開始一共呼叫 $27$ 次函式。

從過程可以看出規律：$A(1, n) = A(1, n - 1) + 1$，$A(1, 0) = 2$，所以 $A(1, n) = n + 2$；
$A(2, n) = A(1, A(2, n - 1)) = A(2, n - 1) + 2$，$A(2, 0) = 3$，所以 $A(2, n) = 2n + 3$，代
$n = 2$ 也得 $7$。 再往上 $A(3, n) = 2^{n+3} - 3$，$A(4, 2) = 2^{65536} - 3$ 已經有將近兩萬位數，
這就是題目說「$m$、$n$ 很小時成長也相當快」。
\end{pastsoln}
